\section{Marco Teórico}

El Marco Teórico de esta investigación es la teoría económica marxista.

\subsection{Perspectiva Crítica: Materialismo Histórico y Economía Marxista}
FALTA DESARROLLAR CON BIBLIOGRAFÍA DE TEORÍA...PENDIENTE...

El análisis se fundamenta en el materialismo histórico y la teoría económica marxista, concibiendo al salario monetario como la expresión del costo de reproducción de la fuerza de trabajo, pagado por el empleador como parte de su capital variable. Bajo este enfoque crítico, la brecha entre los ingresos del hogar y el costo de vida evidencia las contradicciones intrínsecas de la acumulación de capital, la extracción de plusvalía y el deterioro de las condiciones de la clase trabajadora en las economías periféricas y dependientes de Chile, Brasil y Argentina durante el periodo 2015--2025.
